% from ega4-18-fr.tex
\oldpage[IV-4]{185}

(II, 3.7.9, (ii)). En vertu de (18.12.16), $\mathscr L^{\otimes n}$ est donc ample relativement à $Z$ (ce qui équivaut à dire qu'il est ample (II, 4.6.6)), et par suite il en est de même de $\mathscr L$ (II, 4.5.6, (i)).

Nous laissons au lecteur le soin d'énoncer de même des conditions équivalentes à celles de (II, 4.6.3) pour les $\sh{O}_X$-Modules relativement amples.

\section{Immersions régulières et platitude normale}
\label{section:IV.19}

Le présent paragraphe est consacré, d'une part à une étude des immersions régulières (16.9.2) $Y\to X$ de préschémas, notamment dans le cas où $X$ et $Y$ sont plats sur un même préschéma $S$; d'autre part, il donne des compléments sur les suites $M$-régulières et la platitude normale, en généralisant notamment des résultats de platitude établis par H. Hironaka dans sa théorie de la résolution des singularités [35].

\subsection{Propriétés des immersions régulières.}
\EGAReaderTarget{ega4.c04.subsec.19-1}{subsection:IV.19.1}

\begin{proposition}[19.1.1]
\label{IV.19.1.1}
Soient $X$ un préschéma localement noethérien, $j:Y\to X$ une immersion, $y$ un point régulier de $Y$. Pour que $j$ soit une immersion régulière au point $y$, il faut et il suffit que $X$ soit régulier au point $j(y)$.
\end{proposition}

Compte tenu de (16.9.10), on est ramené à prouver le

\begin{corollary}[19.1.2]
\EGAReaderTarget{ega4.c04.cor.19-1-2}{IV.19.1.2}
Soient $A$ un anneau local noethérien, $\mathfrak J$ un idéal de $A$; on suppose que l'anneau quotient $B=A/\mathfrak J$ soit régulier. Pour que $A$ soit régulier, il faut et il suffit que l'idéal $\mathfrak J$ soit régulier.
\end{corollary}

En effet, si $A$ est régulier, on sait ($0$, 17.1.9) que $\mathfrak J$ est engendré par une partie d'un système régulier de paramètres de $A$, donc $\mathfrak J$ est régulier. Inversement, si $\mathfrak J$ est régulier, et si $(x_i)_{1\leq i\leq r}$ est une suite $A$-régulière engendrant $\mathfrak J$, $(x_i)$ fait partie d'un système de paramètres de $A$ ($0$, 16.4.1), donc on déduit de ($0$, 17.1.7) que $A$ est régulier.

\begin{definition}[19.1.3]
\EGAReaderTarget{ega4.c04.def.19-1-3}{IV.19.1.3}
Soient $X$ un préschéma, $Y$ un sous-préschéma de $X$, $U$ un ouvert de $X$ contenant $Y$ et tel que $Y$ soit défini par un Idéal $\mathscr J$ de $\sh{O}_U$; on suppose que $\mathscr J/\mathscr J^2$ soit un $(\sh{O}_U/\mathscr J)$-Module localement libre (ce qui est le cas si $Y$ est quasi-régulièrement immergé dans $X$). Pour tout $y\in Y$, on appelle \emph{codimension transversale de $Y$ dans $X$ au point $y$} et on note $\operatorname{codim}_y^*(Y,X)$ le rang du $(\sh{O}_{U,y}/\mathscr J_y)$-module libre $\mathscr J_y/\mathscr J_y^2$. Si $f:Z\to X$ est une immersion d'image $Y$, on appelle de même \emph{codimension transversale de $f$ au point $z\in Z$} la codimension transversale dans $X$ au point $f(z)$ du sous-préschéma $Y$.
\end{definition}

\begin{proposition}[19.1.4]
\label{IV.19.1.4}
Soient $X$ un préschéma localement noethérien, $Y$ un sous-préschéma de $X$ régulièrement immergé; alors, pour tout $y\in Y$, on a
\begin{equation}
  \operatorname{codim}_y^*(Y,X)=\operatorname{codim}_y(Y,X).
  \tag{19.1.4.1}\label{IV.19.1.4.1}
\end{equation}
\end{proposition}

En vertu de (5.1.3.2), $\operatorname{codim}_y(Y,X)$ est égal à la plus petite valeur des dimensions des anneaux locaux $A=\sh{O}_{X,z}$, où $z$ est point générique d'une composante

